\section{Description entropy under repeated adaptive use}\label{sec:adaptiveentropy67}
The number of calls changes the precision needed in a reusable
instrument description. On a strictly positive Choi body the relevant
scale is $N^{-1/2}$, not the $N^{-1}$ scale supplied by a general
one-slot hybrid. We prove a finite, multiparameter coding statement
against arbitrary common adaptive quantum testers.

For two instruments $J,L\in\mathcal C_{d,n,m}$, define
\begin{equation}\label{eq:adaptivemetric67}
 d_N(J,L)=\sup_{\mathsf T}
       \|\Omega_J^{\mathsf T}-\Omega_L^{\mathsf T}\|_1.
\end{equation}
Here $\mathsf T$ is the same causal tester in both experiments, uses the
same memoryless instrument at most $N$ times, and may have arbitrary
finite-dimensional quantum memory, initially entangled references,
intermediate quantum channels and measurements, and classical
feedback. Outcomes and final memory are retained. A public stopping
time bounded by $N$ is allowed. A changed input dimension can be handled
by a tester channel between calls. There is no coherent superposition
of command ports, no tester informed of the hidden choice $J$ versus
$L$, and no stopping based on inaccessible simulator state. Taking
suprema of trace distances gives the triangle inequality. In
particular $d_N\le2$. Public stopping is included by an absorbing flag,
as in Theorem~\ref{thm:adaptivediamond66}.

For a fixed $a$ with $0<a<1/(mn)$ let
\[
 \mathcal C_a=\{J\in\mathcal C:J_y\ge aI_{dn}\text{ for every }y\}.
\]
This is a Choi eigenvalue margin, not a spectral gap for an action or a
mixing hypothesis on trajectories. Its affine dimension is the same
$s=d^2(mn^2-1)$ as before.

\begin{lemma}[Two-point classical programme bound]\label{lem:programme67}
Suppose $J_y,L_y\ge\lambda I_{dn}$ for every outcome and
$r=|J-L|_2\le\lambda$, where $\lambda>0$. Then, for every $N\ge1$,
\begin{equation}\label{eq:rootN67}
 d_N(J,L)\le\min\{2,\,2\sqrt N\,r/\lambda\}.
\end{equation}
The same bound covers arbitrary initially entangled references and
bounded public stopping in \eqref{eq:adaptivemetric67}.
\end{lemma}
\begin{proof}
There is nothing to prove for $r=0$. Put $K=(J+L)/2$ and
$U=(J-L)/r$. Since $U\in V$ and each $\|U_y\|_{\mathrm{op}}\le1$,
\[
                      K^+=K+\lambda U,\qquad K^-=K-\lambda U
\]
are legal instruments. With
\[
 p_\pm=\tfrac12\pm\frac r{4\lambda},\qquad
 q_\pm=\tfrac12\mp\frac r{4\lambda},
\]
we have $J=p_+K^++p_-K^-$ and $L=q_+K^++q_-K^-$.
All four probabilities are at least $1/4$. The elementary inequality
$\log x\le x-1$ implies
\[
 D(p\|q)\le\sum_{b\in\{+,-\}}\frac{(p_b-q_b)^2}{q_b}
                                      \le\frac{2r^2}{\lambda^2}.
\]
Choose the programme signs independently in advance for all $N$
slots. Conditional on a sign string, a fixed tester acts on exactly
the same sequence of genuine quantum instruments in the two
experiments. It therefore defines one channel from the classical sign
string to the final quantum--classical state. This remains true with
adaptive feedback: conditioning fixes the instruments, not the
tester's outcomes. If the tester stops early, the unused signs are
ignored. Contractivity and the classical product identity give
\[
 \|\Omega_J^{\mathsf T}-\Omega_L^{\mathsf T}\|_1
 \le2\operatorname{TV}(p^{\otimes N},q^{\otimes N})
 \le\sqrt{2N D(p\|q)}\le2\sqrt N\,r/\lambda.
\]
The middle bound is the usual Pinsker inequality. To recall an elementary
proof, coarse-grain two distributions to the set where the first is
larger. The log-sum inequality reduces relative entropy to binary
relative entropy. Its second derivative in the first argument is
$1/[x(1-x)]\ge4$, so integration about its minimum proves
$D\ge2\operatorname{TV}^2$. Taking the supremum over testers proves
the assertion.
\end{proof}
The classical-programme idea is an established method in quantum
metrology; Demkowicz-Dobrza\'nski, Ko\l ody\'nski and Gu\c{t}\u{a}
\cite[Equations~(4)--(6), Methods]{DKG67} use parameter-dependent
mixtures of fixed channels to bound distinguishability. The lemma
above gives the finite two-point trace-norm estimate needed here,
including common adaptive testers and public stopping. It is not a
claim that the mixture method or the square-root scale is a new
metrological principle. The operations $K^\pm$ remain quantum
operations, not classical physical substitutes for unknown input.

\begin{lemma}[Repeated Choi tests]\label{lem:choitesting67}
For any two instruments $J,L\in\mathcal C$, with $t=|J-L|_2$,
\begin{equation}\label{eq:choitesting67}
 d_N(J,L)\ge2-4\exp\left(-\frac{Nt^2}{8d^2}\right).
\end{equation}
A negative right-hand side is simply a vacuous lower bound.
\end{lemma}
\begin{proof}
At each call supply half of a normalized maximally entangled state and
retain its other half and the outcome. The resulting density matrices
are $A=d^{-1}\bigoplus_yJ_y$ and $B=d^{-1}\bigoplus_yL_y$.
They have trace one. The projector on the positive part of $A-B$
gives a binary measurement whose probability gap is
\[
 g=\tfrac12\|A-B\|_1\ge t/(2d).
\]
The same binary measurement is made on each of the $N$ independent
outputs. With a threshold halfway between the two means, each wrong-side
probability is at most $\exp(-Ng^2/2)$. Indeed for a Bernoulli variable
$Z$, the second derivative of
$\log\mathbb E e^{u(Z-\mathbb EZ)}$ is the variance of a tilted
Bernoulli variable and is at most $1/4$. The value and first derivative
at zero vanish; integration and Chernoff's inequality give
$\Pr\{N^{-1}\sum Z_i-\mathbb EZ\ge v\}\le e^{-2Nv^2}$,
and likewise for the other tail. Set $v=g/2$.
The total variation distance between the two product experiments is
at least $1-2e^{-Ng^2/2}$; the unhalved trace norm of their classical
outputs is twice this. This particular nonadaptive tester is allowed
in \eqref{eq:adaptivemetric67}, and the result follows. The tester may
depend on the pair $J,L$, as is permitted when lower-bounding a metric;
it is the same tester for the two members of that pair.
\end{proof}

\begin{theorem}[Sharp horizon exponent for interior descriptions]\label{thm:adaptiveentropy67}
Fix $d,n,m$ with $s>0$, a margin $0<a<1/(mn)$ and a tolerance
$0<\delta<1$. Let $\mathcal M_N(\delta;a)$ be the minimum number of
legal decoded instruments whose $d_N$-balls of radius $\delta$ cover
$\mathcal C_a$. A decoded instrument need not itself have margin $a$.
Write $\mathcal M_N^{\mathbb Q}$ when rational decoded instruments are
required. Set
\[
 R=\frac1{mn}-a,\qquad
 c_\delta=\sqrt{8d^2\log\frac4{1-\delta}},\qquad
 B_N=\left\lceil\frac Ka\left(2+
                         \frac{4\lceil\sqrt N\rceil}{\delta}\right)\right\rceil,
\]
where $K$ is from \eqref{eq:codeconstants67}. Then
\begin{equation}\label{eq:adaptivecover67}
 \max\{1,(R\sqrt N/c_\delta)^s\}
 \le\mathcal M_N(\delta;a)
 \le\mathcal M_N^{\mathbb Q}(\delta;a)
 \le(2B_N+1)^s.
\end{equation}
Consequently the minimum fixed-length description satisfies
\begin{equation}\label{eq:adaptivebits67}
 \log_2\mathcal M_N^{\mathbb Q}(\delta;a)
       =\frac s2\log_2N+O_{d,n,m,a,\delta}(1),
\end{equation}
and the same formula holds without the rationality restriction.
The upper bound is achieved by the exact intrinsic codec applied once
at grid $B_N$; its decoded instrument is reused at every call.
\end{theorem}
\begin{proof}
The tuple-Frobenius ball in $V$ of radius $R$ about $J^\circ$ is
contained in $\mathcal C_a$. Choose a maximal set of points in that
ball with pairwise distances greater than $c_\delta/\sqrt N$.
Balls of that radius about the selected points cover the original
ball. Euclidean volume comparison therefore gives at least
$(R\sqrt N/c_\delta)^s$ points. By Lemma~\ref{lem:choitesting67}, their
pairwise $d_N$ distances are greater than
\[
 2-4e^{-c_\delta^2/(8d^2)}=1+\delta>2\delta.
\]
One $d_N$-ball of radius $\delta$ cannot cover two of these points,
regardless of the location of its legal centre. This proves the lower
bound.

For the upper bound use Theorem~\ref{thm:intrinsiccodec67}. It gives
\[
 r=|J-\widehat J|_2\le\sum_y\|J_y-\widehat J_y\|_1
 \le K/B_N\le\min\{a/2,\ a\delta/(4\sqrt N)\}.
\]
Every block of $\widehat J$ is consequently at least $(a/2)I$, and
both $J$ and $\widehat J$ satisfy Lemma~\ref{lem:programme67} with
$\lambda=a/2$ and $r\le\lambda$. Its conclusion gives
$d_N(J,\widehat J)\le4\sqrt N\,r/a\le\delta$.
There are at most $(2B_N+1)^s$ possible digit strings. All codewords
arising from targets in $\mathcal C_a$ decode legally. The remaining
codewords may be rejected or assigned $J^\circ$; this does not affect
the cover. Since $B_N=O_{d,n,m,a,\delta}(\sqrt N)$, taking logarithms
in \eqref{eq:adaptivecover67} proves \eqref{eq:adaptivebits67}.
\end{proof}

\begin{remark}[What is charged]\label{rem:adaptivescope67}
The lower bound concerns the number of possible reusable mathematical
instrument descriptions, not the mutable memory of a simulator that
receives and stores the entire instrument as input. It applies to all
full-dimensional strictly positive instrument families above, without
an expanding unitary subsystem. It therefore complements, rather than
replaces, the hypothesis-sensitive streaming-space lower bounds.
For variable $a,\delta,d,n,m$, the displayed finite inequalities remain
valid; \eqref{eq:adaptivebits67} holds with those parameters fixed.
We assert neither an optimal uniform dimension constant nor a complete
joint $N$--$\delta$ asymptotic. The rational implementation takes
rational $a,\delta$ and an asserted margin, which it can verify exactly.
No samples of an unknown channel are supplied or required.
\end{remark}

\begin{proposition}[The boundary is quantitatively different]\label{prop:boundaryphase67}
There is no constant $C$ such that
$d_N(J,L)\le C\sqrt N\,|J-L|_2$ holds for every pair of qubit channels
and every $N$. In particular the strict margin in
Theorem~\ref{thm:adaptiveentropy67} cannot be removed from its proof by
setting $a=0$.
\end{proposition}
\begin{proof}
Let $U_\theta=\operatorname{diag}(1,e^{i\theta})$ and compare its
unitary channel to the identity channel. For their unnormalized Choi
matrices,
\[
 |J(U_\theta)-J(I)|_2
       =2\sqrt2\,|\sin(\theta/2)|\le\sqrt2\,|\theta|.
\]
At $\theta=\pi/N$, sequentially applying the unknown channel $N$
times to $|+\rangle$ gives respectively $|-\rangle$ and $|+\rangle$.
The final trace distance is two. The proposed upper bound would tend
to zero. The general additive diamond hybrid remains valid at this
boundary; the proposition asserts no global covering exponent for
boundary or rank-constrained instrument classes.
\end{proof}

\subsection{Description, learning and physical simulation}
A reusable rational description is different from a protocol that
replaces quantum communication by classical messages.
Naik, Zartab, Gisin and Banik \cite[Theorems~2 and~5]{NZGB67} show a
finite-classical-message obstruction for the perfect qubit channel in
a reference-sensitive joint-measurement model, and finite-resource
possibility for noisy depolarizing channels. Their receiver retains a
quantum system; the task is physical simulation. Here the decoded
objects and the vertices in Lemma~\ref{lem:programme67} are genuine
quantum instruments. Our numerical algorithms instead start from
specified matrices. Neither construction supplies a classical physical
device accepting unknown entangled input.

Unknown-channel learning charges channel queries used to infer a
classical description. Mele and Bittel \cite[Theorem~I.1]{MB67}
prove dimension-optimal query complexity at fixed diamond accuracy;
Oufkir and Girardi \cite{OG67} give further accuracy-dependent query
lower bounds. The input here is already specified, and the lower bound
in \eqref{eq:adaptivecover67} counts descriptions, not the queries
needed to select one from unknown data. The elementary packing and
testing ingredients are shared methods, not a priority claim that
those learning theorems are consequences of this code.
